

\section{Preliminaries and Problem Formulation} \label{sec:pb}


This section presents the notation, set representations, system dynamics, and reachable set definitions used in this work. We then pose our safe RL problem.

\subsection{Notation and Set Representations}\label{subsec:notation}
%\textbf{Notation and Set Representations.}
The $n$-dimensional real numbers are $\R^n$, the natural numbers are $\N$, and the integers from $n$ to $m$ are $n{:}m$.
% Scalars are lowercase italic and sets are uppercase italic.
% Vectors are lowercase bold and matrices are uppercase bold.
% Set-valued functions are in uppercase script.
% For vectors, we denote the element $i$ of vector $\vc{v}$ by $\vc{v}\arridx{i}$. 
% We denote the element at row $i$ and column $j$ of matrix $A$ by $A\arridx{i,j}$.
% We denote column $j$ of $A$ by $A\arridx{\cdot,j}$.
We denote the element at row $i$ and column $j$ of matrix $\vc{A}$ by %$\arridx{\vc{A}}{i,j}$.
$(\vc{A})_{i,j}$, column $j$ of $\vc{A}$ by $\arridx{\vc{A}}{:\,,j}$, and the element $i$ of vector $\vc{a}$ by $\arridx{\vc{a}}{i}$.
An $n\times m$ matrix of ones is $\ones_{n\times m}$.
For $\vc{A} \in \R^{n\times m}$ and $\vc{x} \in \R^n$, we use the shorthand $\vc{A} - \vc{x} = \vc{A} - \vc{x}\ones_{1\times m}$.
The $\diag{\cdot}$ operator places its arguments block-diagonally in a matrix of zeros.
% The Kronecker product is denoted $\otimes$.
For a pair of sets $A$ and $B$, the Minkowski sum is $A + B = \{\vc{a}+\vc{b}\ |\ \vc{a} \in A, \vc{b} \in B\}$, and the Cartesian product is $A \times B = \{(\vc{a},\vc{b})\ |\ \vc{a} \in A, \vc{b} \in B\}$.
%The $i$\ts{th} element of $A$ is $a\idxi \in A$. 



% \subsection{Set Representations}
We represent sets using constrained zonotopes, zonotopes, and intervals, because they enable efficient Minkowski sum computation (a key part of reachability analysis) \cite{althoffphdthesis} and collision checking via linear programming (critical to safe motion planning) \cite{conf:const_zono}.
A \emph{constrained zonotope} \cite{conf:const_zono} is a convex set parameterized by a center $\ctr \in \R^n$, generator matrix $\Gen \in \R^{n \times \ngen}$, constraint matrix $\Acon \in $ $\R^{\ncon \times \ngen}$, and constraint vector $\bcon \in \R^{\ncon}$ as
\begin{equation}\label{eq:conszono}
    \zono{\ctr,\Gen,\Acon,\bcon} = \big\{
        \ctr + \Gen\coef\ |\ \Acon\coef = \bcon,\ \norm{\coef}_\infty \leq 1
    \big\}.
\end{equation}
By \cite[Thm. 1]{conf:const_zono}, every convex, compact polytope is a constrained zonotope and vice-versa.
For polytopes represented as an intersection of halfplanes, we convert them to constrained zonotopes by finding a bounding box, then applying the halfspace intersection property in \cite{raghuraman2020set_ops_conzono}.

A zonotope is a special case of a constrained zonotope without equality constraints (but with $\norm{\coef}_\infty \leq 1$), which we denote $\zono{\ctr,\Gen}$.
For $Z = \zono{\ctr,\Gen} \subset \R^n$ and a linear map $L$, we have $LZ  = \zono{L\ctr, L\Gen}$; we denote $-Z = -1Z$.
The Minkowski sum of two zonotopes $Z_1 = \zono{\ctr_1,\Gen_1}$ and $Z_2 = \zono{\ctr_2,\Gen_2}$ is given by $Z_1 + Z_2 = \zono{\ctr_1+\ctr_2,[\Gen_1,\Gen_2]}$ \cite{althoffphdthesis}.
For an $n$-dimensional interval with lower (resp. upper) bounds $\lb \in \R^n$ (resp. $\ub$), we abuse notation to represent it as a zonotope $Z = \zono{\lb,\ub} \subset \R^n$, with center $\tfrac{1}{2}(\lb+\ub)$ and generator matrix $\diag{\tfrac{1}{2}(\ub - \lb)}$.



\subsection{Robot and Environment}
%\textbf{Robot and Environment.}
We assume the robot can be described as a discrete-time, nonlinear control system with state $\vcstate_{k} \in \statespace \subset \R^n$ at time $k \in \N$.
We assume the state space $\statespace$ is compact.
The input $\action_{k}$ is drawn from a zonotope $\actionspace_k \subseteq \actionspace$ at each time $k$, where $\actionspace \subset \R^m$ is a zonotope of all possible actions.
We denote process noise by $\noise_{k} \in \noisespace \subset \R^n$, where $\noisespace$ is specified later in Assumption \ref{assumption:noise_zonotope}.
Finally, we denote the black box (i.e., unknown) dynamics $\dyn: \statespace\times\actionspace\times\noisespace \to \statespace$, for which
\begin{align}\label{eq:sys}
        \vcstate_{k+1} &=  \dyn(\vcstate_{k}, \action_{k}) + \noise_{k}.
\end{align}
We further assume that $\dyn$ is twice differentiable and Lipschitz continuous, meaning there exists a \emph{Lipschitz constant} $L^\star$ such that, if $\forall$ $\vc{z}_1, \vc{z}_2 \in \R^{n+m}$ with $\vc{z}_j= (\vcstate_j,\action_j)$, then $\norm{\dyn(\vc{z}_1) - \dyn(\vc{z}_2)} \leq L^\star \norm{\vcstate_1 - \vcstate_2}$.
We denote the initial state of the system as $\vcstate_{0}$, drawn from a compact set $\statespace_{0} \subset \R^n$.
Note that this formulation leads to an MDP.

To enable safety guarantees, we leverage the notion of failsafe maneuvers from mobile robotics \cite{kousik2020bridging,magdici2016fail}. % as follows:
\begin{assumption}\label{assumption:stop}
We assume the dynamics $\dyn$ are invariant to translation in position, and the robot can brake to a stop in $\nbrk \in \N$ time steps and stay stopped indefinitely.
That is, there exists $\action\brk \in \actionspace$ such that, if the robot is stopped at state $\vcstate_k$, and if $\vcstate_{k+1} = \dyn(\vcstate_k,\action\brk)$, then $\vcstate_{k+1} = \vcstate_k$.
\end{assumption}
\noindent Note, many real robots have a braking safety controller available, similar to the notion of an invariant set \cite{LewEtAl2021,conf:modelBasedReachRL}.
Also, failsafe maneuvers exist even when a robot cannot remain stationary, such loiter circles for aircraft \cite{fridovich2019safely,althoff2015online}.
% Since the dynamics $\dyn$ are Lipschitz, and $\statespace$ is compact (velocity is bounded), there exists a finite \emph{braking distance}, denoted $d\brk > 0$.
%We investigate the impact of this assumption on our method's performance with a hexacopter in wind numerical example in Section \ref{sec:eval}.

We require that process noise obeys the following assumption for numerical tractability and robustness guarantees.

\begin{assumption}\label{assumption:noise_zonotope}
Each $\noise_{k}$ is drawn uniformly from a \emph{noise zonotope} $\noisespace = \zono{\ctr_\noise,\Gen_\noise}$ with $\nnoisegen$ generators.% (which is not time-varying).
\end{assumption}
\noindent This formulation does not handle discontinuous changes in noise.
However, there exist zonotope-based techniques to identify a change in $\noisespace$ \cite{shetty2020predicting}, after which one can compute the system's reachable set as in the present work.
We leave measurement noise and perception uncertainty to future work.
We also note, in the case of Gaussian or unbounded noise, one can overapproximate a confidence level set of a probability distribution using a zonotope \cite{althoffphdthesis,shetty2020predicting}.
% Indeed, safe navigation with measurement noise or inexact state information is an active area of research that we leave to future work.

We denote unsafe regions of state space, or \emph{obstacles}, as $\statespace\obs \subset \statespace$.
We assume obstacles are static but different in each episode, as the focus of this work is not on predicting other agents' motion.
Furthermore, reachability-based frameworks exist to handle other agents' motion \cite{leung2020infusing,vaskov2019not}, so the present work can extend to dynamic environments.

We further assume the robot can instantaneously sense all obstacles (that is, $\statespace\obs$) and represent them as a union of constrained zonotopes.
%We also note that common obstacle representations, such as occupancy grids, are polygonal, so they can be represented as constrained zonotopes.
In the case of sensing limits, one can determine a minimum distance within which obstacles must be detected to ensure safety, given a robot's maximum speed and braking distance \cite[Section 5]{kousik2020bridging}.





\subsection{Reachable Sets}\label{subsec:reachable_sets_defn}
%\textbf{Reachable Sets.}
We ensure safety by computing our robot's forward reachable set (FRS) for a given motion plan, then adjusting the plan so that the FRS lies outside of obstacles.
We define the FRS, henceforth called the reachable set, as follows:
\begin{definition} 
The reachable set $\reachset_{k}$ at time step $k$, subject to a sequence of inputs $\action_{j} \in \actionspace_{j} \subset \mathbb{R}^m$, noise $\noise_{j} \in \noisespace$ $\forall\ j \in \{ 0, \dots, k-1\}$, and initial set $\statespace_{0} \in  \mathbb{R}^n$, is the set
\begin{align}\begin{split}\label{eq:reachable_set_k}
        \reachset_{k} = \big\{&\vcstate_{k} \in \mathbb{R}^n \, \big|\ 
            \vcstate_{j+1} = \dyn(\vcstate_{j}, \action_{j}) + \noise_{j},\ \vcstate_{0} \in \statespace_{0},\\ &\action_{j} \in \actionspace_{j},\ \regtext{and}\ \noise_{j} \in \noisespace,\  \forall\ j = 0,\cdots,k-1\big\}.
\end{split}\end{align}
\end{definition}

Recall that we treat the dynamics $\dyn$ as a black box (e.g., a simulator), which could be nonlinear and difficult to model, but we still seek to conservatively approximate (that is, overapproximate) the reachable set $\reachset_k$.
% Doing so requires conservatively estimating the Lipschitz constant of the dynamics, as is done in the literature \cite{koller2018learning,alanwar2020data_conf}.

% By computing a conservative reachable set, we can conservatively approximate the Lipschitz constant of the dynamics.
% Consider the set $F = \cup_{j=0}^{k} (\reachset_{j} \times \actionspace_{j})$.
% We can approximate $L^\star \geq 0$ such that $\| \dyn(\vc{y}_1) - \dyn(\vc{y}_2) \|_2 \leq L^\star \| \vc{y}_1 - \vc{y}_2\|_2$ holds for all $\vc{y}_1, \vc{y}_2 \in F$.



\subsection{Safe RL Problem Formulation}
%\textbf{Problem Formulation.}
We denote the state of the RL agent at time $k$ by  $\RLstate_{k} \in \R^\nrl$, which contains the state $\vcstate_k$ of the robot plus information such as sensor measurements and previous actions.
At each time $k$, the RL agent chooses $\RLaction_{k}$.
Recall that $\statespace\obs \subset \R^n$ denotes obstacles.
For a given task, we construct a reward function $\rewfunc: (\RLstate_{k},\action_k) \mapsto \rew_{k} \in \R$ (examples of $\rewfunc$ are given in Section \ref{sec:eval}).
At time $k$, let $\plan_k = (\action_j)_{j=k}^\nplan$ denote a \emph{plan}, or sequence of actions, of duration $\nplan \in \N$.

Then, our safe RL problem is as follows.
We seek to learn a policy $\policy_\theta: \RLstate_{k} \mapsto \action_{k}$, represented by a neural network with parameters $\theta$, that maximizes expected cumulative reward.
Note that the policy can be deterministic or stochastic.
Since rolling out the policy na\"ively may lead to collisions, we also seek to create a safety layer between the policy and the robot (that is, to ensure $\reachset_j \cap \statespace\obs = \emptyset$ for all $j \geq k$).

% Next, we detail our proposed safe RL method.